% ============================================================
% 60_boxes.tex -- Balken, generische Boxen und Listen
% ============================================================

% ------------------------------------------------------------
% Kapitel- und Abschnittsbalken
% ------------------------------------------------------------
\newtcolorbox{chapterbar}[1][]{
  enhanced jigsaw,
  colback=\ChapterBarColor,
  colframe=\ChapterBarColor,
  boxrule=0pt,
  boxsep=0pt,
  arc=\ZSFcornerBox, outer arc=\ZSFcornerBox,
  left=\ZSFspaceS, right=\ZSFspaceS,
  top=\ZSFspaceS, bottom=\ZSFspaceS,
  before skip=0pt,
  after skip=0pt,
  fontupper=\ZSFfontChapter,
  #1
}
\newcommand{\ChapterBar}[1]{%
  \ZSFHeadingBefore{\ZSFchapterBeforeSkip}%
  \begin{chapterbar}\textcolor{\ChapterBarTextColor}{#1}\end{chapterbar}%
  \ZSFHeadingAfter{\ZSFchapterAfterSkip}%
}

\newtcolorbox{subsectionbar}[1][]{
  enhanced jigsaw,
  colback=\SubsectionBarColor,
  colframe=\SubsectionBarColor,
  boxrule=0pt,
  boxsep=0pt,
  arc=\ZSFcornerBox, outer arc=\ZSFcornerBox,
  left=\ZSFspaceS, right=\ZSFspaceS,
  top=.75\ZSFspaceXS, bottom=.75\ZSFspaceXS,
  before skip=0pt,
  after skip=0pt,
  fontupper=\ZSFfontSection,
  #1
}
\newtcolorbox{subsubsectionbar}[1][]{
  enhanced jigsaw,
  colback=\SubsubsectionBarColor,
  colframe=\SubsubsectionBarColor,
  boxrule=0pt,
  boxsep=0pt,
  arc=\ZSFcornerBox, outer arc=\ZSFcornerBox,
  left=\ZSFspaceS, right=\ZSFspaceS,
  top=\ZSFspaceXS, bottom=\ZSFspaceXS,
  before skip=0pt,
  after skip=0pt,
  fontupper=\ZSFfontSubsubsection,
  #1
}

\makeatletter
\newcommand{\SubsectionBar}{\@ifstar{\SubsectionBarStar}{\SubsectionBarNumbered}}
\newcommand{\SubsectionBarNumbered}[2][]{%
  \refstepcounter{subsection}%
  \setcounter{subsubsection}{0}%
  \if\relax\detokenize{#1}\relax\else\label{#1}\fi
  \ZSFHeadingBefore{\ZSFsubsectionBeforeSkip}%
  \begin{subsectionbar}\textcolor{\SubsectionBarTextColor}{\thesubsection. #2}\end{subsectionbar}%
  \ZSFHeadingAfter{\ZSFsubsectionAfterSkip}%
}
\newcommand{\SubsectionBarStar}[2][]{%
  \if\relax\detokenize{#1}\relax\else\label{#1}\fi
  \ZSFHeadingBefore{\ZSFsubsectionBeforeSkip}%
  \begin{subsectionbar}\textcolor{\SubsectionBarTextColor}{#2}\end{subsectionbar}%
  \ZSFHeadingAfter{\ZSFsubsectionAfterSkip}%
}
\newcommand{\SubsubsectionBar}{\@ifstar{\SubsubsectionBarStar}{\SubsubsectionBarNumbered}}
\newcommand{\SubsubsectionBarNumbered}[2][]{%
  \refstepcounter{subsubsection}%
  \if\relax\detokenize{#1}\relax\else\label{#1}\fi
  \ZSFHeadingBefore{\ZSFsubsubsectionBeforeSkip}%
  \begin{subsubsectionbar}\textcolor{\SubsubsectionTitleColor}{\thesubsubsection. #2}\end{subsubsectionbar}%
  \ZSFHeadingAfter{\ZSFsubsubsectionAfterSkip}%
}
\newcommand{\SubsubsectionBarStar}[2][]{%
  \if\relax\detokenize{#1}\relax\else\label{#1}\fi
  \ZSFHeadingBefore{\ZSFsubsubsectionBeforeSkip}%
  \begin{subsubsectionbar}\textcolor{\SubsubsectionTitleColor}{#2}\end{subsubsectionbar}%
  \ZSFHeadingAfter{\ZSFsubsubsectionAfterSkip}%
}
\makeatother

% ------------------------------------------------------------
% Gemeinsame Box-Stile
% ------------------------------------------------------------
\tcbset{zsftitlebar/.style={
  fonttitle=\ZSFfontBoxTitle\strut,
  halign title=flush left,
  title style={minimum height=\ZSFtitleBarHeight},
  lefttitle=\ZSFspaceS,
  righttitle=\ZSFspaceS,
  toptitle=0pt,
  bottomtitle=0pt,
  before title={},
  after title={},
}}

\newcommand{\ZSFTitleTag}[1]{\hfill{\ZSFfontTitleTag #1}}

% In multicols* verwaltet ausschliesslich multicol die Spaltenumbrueche.
% Inhaltsboxen bleiben unteilbar; lange Inhalte werden semantisch aufgeteilt.
\tcbset{zsfcontentbase/.style={
  enhanced jigsaw,
  colback=\ContentBoxBackColor,
  colframe=\ContentBoxFrameColor,
  boxrule=\ZSFborderThin,
  boxsep=0pt,
  arc=\ZSFcornerBox, outer arc=\ZSFcornerBox,
  before skip=\ZSFboxBeforeSkip,
  after skip=\ZSFboxAfterSkip,
  before upper={\parskip=\ZSFspaceXS\relax\ZSFReadableAuto},
}}

\tcbset{zsfcontenttitled/.style={
  zsfcontentbase,
  colbacktitle=\BoxTitleColor,
  coltitle=\BoxTitleTextColor,
  colframe=\TitledBoxFrameColor,
  titlerule=0pt,
  zsftitlebar,
}}

\tcbset{zsftitlebox/.style={
  zsfcontenttitled,
  colback=\BoxBackColor,
  colframe=\BoxFrameColor,
  boxrule=\ZSFborderBox,
}}

% ------------------------------------------------------------
% Semantische Inhaltsboxen
% ------------------------------------------------------------
\newtcolorbox{zsf@contentbox@plain}[1][]{
  zsfcontentbase,
  left=\ZSFspaceS, right=\ZSFspaceS,
  top=\ZSFspaceXS, bottom=\ZSFspaceXS,
  #1
}
\newtcolorbox{zsf@contentbox@titled}[3][]{
  zsfcontenttitled,
  left=\ZSFspaceS, right=\ZSFspaceS,
  top=\ZSFspaceXS, bottom=\ZSFspaceXS,
  title={\strut #2\if\relax\detokenize{#3}\relax\else\ZSFTitleTag{#3}\fi},
  #1
}
\NewDocumentEnvironment{contentbox}{o O{}}
  {\IfNoValueTF{#1}{\begin{zsf@contentbox@plain}}{\begin{zsf@contentbox@titled}{#1}{#2}}}
  {\IfNoValueTF{#1}{\end{zsf@contentbox@plain}}{\end{zsf@contentbox@titled}}}

\newtcolorbox{zsf@tablebox@plain}[1][]{
  zsfcontentbase,
  left=\ZSFspaceXS, right=\ZSFspaceXS,
  top=\ZSFspaceXS, bottom=\ZSFspaceXS,
  halign=center,
  before upper={},
  #1
}
\newtcolorbox{zsf@tablebox@titled}[3][]{
  zsfcontenttitled,
  left=\ZSFspaceXS, right=\ZSFspaceXS,
  top=\ZSFspaceXS, bottom=\ZSFspaceXS,
  halign=center,
  before upper={},
  title={\strut #2\if\relax\detokenize{#3}\relax\else\ZSFTitleTag{#3}\fi},
  #1
}
\NewDocumentEnvironment{tablebox}{o O{}}
  {\IfNoValueTF{#1}{\begin{zsf@tablebox@plain}}{\begin{zsf@tablebox@titled}{#1}{#2}}}
  {\IfNoValueTF{#1}{\end{zsf@tablebox@plain}}{\end{zsf@tablebox@titled}}}

\newtcolorbox{defbox}[1][]{zsftitlebox,
  left=\ZSFspaceS, right=\ZSFspaceS,
  top=\ZSFspaceS, bottom=\ZSFspaceS,
  title={#1}}

\newenvironment{longformula}{\[\begin{aligned}}{\end{aligned}\]}
\newtcolorbox{formulabox}[1][]{
  enhanced jigsaw,
  colback=\FormulaboxBackColor,
  colframe=\BoxFrameColor,
  boxrule=\ZSFborderBox,
  arc=\ZSFcornerBox, outer arc=\ZSFcornerBox,
  left=\ZSFspaceS, right=\ZSFspaceS,
  top=\ZSFspaceS, bottom=\ZSFspaceS,
  before skip=\ZSFboxBeforeSkip,
  after skip=\ZSFboxAfterSkip,
  before upper={\centering},
  after upper={\par},
  #1
}

% ------------------------------------------------------------
% Aussagen, Verfahren, Listen
% ------------------------------------------------------------
\tcbset{zsfstatementstyle/.style={
  enhanced jigsaw,
  colback=white,
  colframe=white,
  frame hidden,
  borderline west={\ZSFstatementBarWidth}{0pt}{\StatementBarColor},
  boxrule=0pt,
  arc=0pt, outer arc=0pt,
  left=\ZSFstatementLeftPad,
  right=\ZSFspaceS,
  top=\ZSFspaceXS,
  bottom=\ZSFspaceXS,
  before skip=\ZSFboxBeforeSkip,
  after skip=\ZSFboxAfterSkip,
  before upper={\parskip=\ZSFspaceXS\relax\ZSFReadableAuto},
}}

\NewDocumentEnvironment{statementbox}{O{}}
  {\begin{tcolorbox}[zsfstatementstyle]%
   \if\relax\detokenize{#1}\relax\else
     {\ZSFfontBoxTitle\color{\StatementTitleColor}#1}\par\vspace{\ZSFspaceXS}%
   \fi\ignorespaces}
  {\end{tcolorbox}}

\NewDocumentEnvironment{procedure}{O{}}
  {\begin{statementbox}[#1]%
   \begin{enumerate}[leftmargin=\ZSFprocLeftMargin,labelsep=\ZSFprocLabSep,itemsep=\ZSFprocItemSep,topsep=0pt,parsep=0pt,partopsep=0pt]}
  {\end{enumerate}\end{statementbox}}
\newcommand{\ProcStep}[2]{\item {\ZSFfontBoxTitle #1}\\[\ZSFprocStepGap]#2}

\NewDocumentEnvironment{factlist}{}
  {\begin{itemize}[
     leftmargin=\ZSFlistLeftMargin,
     labelsep=\ZSFprocLabSep,
     label={\textcolor{\StatementBarColor}{\textbullet}},
     itemsep=\ZSFprocItemSep,
     topsep=\ZSFspaceXS,
     parsep=0pt,
     partopsep=0pt
   ]}
  {\end{itemize}}
\newcommand{\ZSFFact}[2]{\item \textbf{#1}\,---\,#2}

\newif\ifzsfIndexFirst
\NewDocumentEnvironment{chapterindex}{}
  {\par\begingroup
   {\ZSFfontSection\color{\chaptercolor} Inhaltsübersicht\par}%
   \vspace{\ZSFspaceXS}%
   \RaggedRight
   \ZSFfontTable
   \setlength{\parskip}{0pt}% Keep entries vertically compact
   \zsfIndexFirsttrue
  }
  {\par\vspace{\ZSFspaceXS}\endgroup}
\newcommand{\ZSFChapterIndexGroup}[2]{%
  \ifzsfIndexFirst
    \zsfIndexFirstfalse
  \else
    \par\vspace{\ZSFspaceS}%
  \fi
  \par\noindent
  \hangindent=1.5em\hangafter=1%
  \makebox[1.5em][l]{\ZSFsectionref{#1}}\ZSFkeyword*{#2}%
}
\newcommand{\ZSFChapterIndexEntry}[2]{%
  \par\noindent
  \hspace*{0.8em}%
  \hangindent=3.2em\hangafter=1%
  \makebox[2.4em][l]{\ZSFsectionref{#1}}#2%
}

\NewDocumentEnvironment{propertylist}{O{}}
  {\begin{statementbox}[#1]\begin{itemize}[
     leftmargin=\ZSFlistLeftMargin,
     labelsep=\ZSFprocLabSep,
     itemsep=\ZSFprocItemSep,
     topsep=0pt,
     parsep=0pt,
     partopsep=0pt
   ]}
  {\end{itemize}\end{statementbox}}

% ------------------------------------------------------------
% Inline-Helfer und Titelkopf
% ------------------------------------------------------------
\newtcbox{\ZSFdanger}{colback=ZSFDangerColor, colframe=ZSFDangerColor, coltext=white,
  on line, boxsep=\ZSFspaceXS, left=\ZSFspaceS, right=\ZSFspaceS,
  top=\ZSFspaceXXS, bottom=\ZSFspaceXXS, boxrule=0pt,
  arc=\ZSFcornerInline, fontupper=\ZSFfontDangerUpper}
\newcommand{\eqbox}[1]{\fcolorbox{\BoxFrameColor}{\chaptercolorlight}{\(\displaystyle #1\)}}
\newcommand{\hl}[1]{\underline{\textbf{#1}}}
\renewcommand{\ZSFhl}[1]{\hl{#1}}

\newtcolorbox{zsftitleheader}[1][]{
  enhanced jigsaw,
  colback=ZSFDocumentTitleBarColor,
  colframe=ZSFDocumentTitleBarColor,
  boxrule=0pt,
  boxsep=0pt,
  arc=\ZSFcornerBox, outer arc=\ZSFcornerBox,
  left=\ZSFspaceM, right=\ZSFspaceM,
  top=\ZSFspaceS, bottom=\ZSFspaceS,
  before skip=0pt,
  after skip=\ZSFspaceXS,
  halign=center,
  #1
}
\newcommand{\ZSFTitleHeader}[2]{%
  \begin{zsftitleheader}
    {\color{ZSFDocumentTitleTextColor}\ZSFfontTitleHeader #1\par}%
    \vspace{\ZSFspaceXXS}%
    {\color{ZSFDocumentTitleBylineColor}\ZSFfontTitleSubheader von~#2\par}%
  \end{zsftitleheader}
}
